\section{Discussion}
\label{sec:discussion}

\paragraph{Where the PINN did well.} On problems with smooth, compatible data the networks reached accuracies
that are respectable for any method: a relative error of $4\times10^{-5}$ on the Taylor--Green vortex, and
0.27\% of the converged field on the $\Rey=1000$ cavity with soft boundary data, below the discretisation
error of our own $256^2$ finite-difference solution. The result is a single differentiable function of space
(and time) that can be evaluated at $3.9$--$12\times10^6$ points per second, differentiated exactly and sampled at
any resolution, which is what the vortex-tube renderings and the drag and lift integrals of this paper
use. The inverse problem, where sparse velocities must be reconciled with the equations and unknown
coefficients, and the parametric surrogate, where one training run serves a range of Reynolds numbers, are
the settings in which this representation is used as intended (Sections~\ref{sec:res-inverse}
and~\ref{sec:res-deeponet}).

\paragraph{Where a classical solver wins.} For a single forward solve the comparison is not close. Our
finite-difference code, an explicit pseudo-time solver chosen for simplicity rather than speed, produced the
$256^2$ regularised-lid solution at $\Rey=1000$ in about ten minutes on the CPU and the $512^2$ solution in 57
minutes; the best PINN needed 1.6~GPU-hours for comparable accuracy and several configurations of the same
method missed by 20--28\%. A Newton or multigrid solver, or a production code such as OpenFOAM or FEniCSx,
would be faster again; we did not run one, so no measured comparison with those codes is claimed. The
unsteady problems are more expensive still: the cylinder and the three-dimensional vortex were given a small
fraction of their planned budgets and still dominate the GPU time of this study.

\paragraph{What the ablation says about the components.} Components that change the \emph{representation}
(Fourier features, the gated network, exact periodicity, exact divergence through a stream function) behaved
as expected wherever they were compatible with the data. Components that change the \emph{objective} were
problem dependent. Exact Dirichlet data are only as good as the data: the regularised lid is continuous but
not compatible with continuity at the corners, and imposing it exactly caused the largest errors in this
study. Gradient-norm balancing, which equalises the gradient contributions of the loss terms, is blind to this:
it amplified a term that is easy to satisfy (the lid) and reduced the one that could not be satisfied, and in
doing so recreated the incompatibility we had removed. The second-order L-BFGS stage, finally, is only as good
as its fixed batch.

\paragraph{The reference data matter as much as the method.} Two of our gates turned out to depend on the
reference more than on the network. The Ghia tables are an excellent 1982 multigrid solution, but at the
2\% level their $v$ profile differs from grid-converged solutions of both the unit-lid and the
regularised-lid problems. For the three-dimensional vortex no numeric reference curve was available with the
workshop descriptions, and we had to digitise one from a figure. A benchmark gate should be stated against a
reference whose own error is well below the gate, and the reference should solve the same problem (here: the
same lid).

\paragraph{Limitations.} All runs use one random seed, so no variance across seeds is reported. The budgets
of the cylinder, three-dimensional, inverse, parametric and ablation runs are reduced (Table~\ref{tab:budget})
to fit the available GPU time; results for those problems are therefore lower bounds on what the method can
do with the planned budgets, and the cylinder ablation was not run at all. For most runs the GPU was shared with
an unrelated workload, so absolute timings are pessimistic by a factor of about 1.3--1.5. The
$64^4$ collocation grid suggested for the separable network did not fit in 16~GB, and the SOAP optimiser was
not used. The three-dimensional reference is digitised from a published figure (about 1\% of full scale). The
drag, lift and Strouhal reference intervals were taken from the FEATFLOW pages because the original benchmark
paper could not be retrieved. No OpenFOAM or FEniCSx ground truth was generated; the cavity reference is our own
finite-difference solution, validated by grid refinement and against the spectral value of the primary vortex.
